\documentclass[11pt,a4paper]{article}
\usepackage[margin=1in]{geometry}
\usepackage{amsmath,amssymb,amsthm}
\usepackage{algorithm}
\usepackage{algpseudocode}
\usepackage{booktabs}
\usepackage{hyperref}

\theoremstyle{definition}
\newtheorem{definition}{Definition}[section]
\newtheorem{theorem}{Theorem}[section]
\newtheorem{lemma}[theorem]{Lemma}
\newtheorem{remark}{Remark}[section]

\title{\textbf{Theoretical Foundations, Ensemble Approximations, and Inverted Scaling Mechanics of Dropout Regularization}}
\author{Team Scaibu \\ \small Email: \texttt{jaydeep.9664920749@gmail.com} \\ \small Architecture and Core Engine Team}
\date{October 2026}

\begin{document}
\maketitle

\begin{abstract}
This document provides the rigorous mathematical specification, expectation-invariance theory, and ensemble interpretation of Inverted Dropout. In standard dropout, random sub-networks are sampled during training, requiring post-hoc weight scaling during evaluation. Inverted dropout scales kept activations by $\frac{1}{1-p}$ during the forward training pass, preserving expectation identically and rendering inference a pure identity passthrough. We formalize the Bernoulli masking process, prove unbiased gradient expectation, detail computational complexity, trace a complete worked example, and analyze edge-case considerations.
\end{abstract}

\section{Notation and Mathematical Preliminaries}

Let $X \in \mathbb{R}^{B \times D}$ denote an intermediate layer activation tensor across batch size $B$ and hidden feature dimension $D$.

\begin{itemize}
    \item $B \in \mathbb{N}_{\ge 1}$: Batch size.
    \item $D \in \mathbb{N}_{\ge 1}$: Activation channel dimension.
    \item $p \in [0, 1)$: Dropout probability (rate of zeroing out units).
    \item $q = 1 - p \in (0, 1]$: Retention probability (survival probability of a unit).
    \item $M_{ij} \sim \text{Bernoulli}(q)$: Independent random binary mask variable for token $i$ and feature $j$.
    \item $s = \frac{1}{q} = \frac{1}{1-p}$: Inverted scaling multiplier during training.
    \item $Y \in \mathbb{R}^{B \times D}$: Output tensor after stochastic masking and scaling.
\end{itemize}

\begin{table}[h]
\centering
\caption{Summary of Mathematical Symbols for Inverted Dropout}
\begin{tabular}{lll}
\toprule
\textbf{Symbol} & \textbf{Type / Domain} & \textbf{Description} \\
\midrule
$X$ & $\mathbb{R}^{B \times D}$ & Input activation matrix \\
$p$ & $[0, 1)$ & Dropout zeroing probability \\
$q$ & $(0, 1]$ & Keep probability ($1 - p$) \\
$M$ & $\{0, 1\}^{B \times D}$ & Stochastic Bernoulli mask matrix \\
$s$ & $[1, \infty)$ & Inverted scale factor ($1/q$) \\
$Y$ & $\mathbb{R}^{B \times D}$ & Regularized output matrix \\
\bottomrule
\end{tabular}
\end{table}

\section{Problem Definition and Core Concepts}

\subsection{Intuitive Meaning}
Deep neural networks often develop co-adaptations, where neurons become overly dependent on the presence of specific correlated features. Dropout breaks these co-adaptations by randomly deactivating a fraction $p$ of units during each training iteration. This forces every unit to learn representations that are independently robust, acting as an implicit geometric ensemble over $2^D$ thinned architectures.

\subsection{Formal Mathematical Definition}
During training ($\text{training} = \text{true}$), each coordinate $(i, j)$ is transformed as:
\begin{align}
M_{ij} &\sim \text{Bernoulli}(1 - p) \\
Y_{ij} &= \frac{1}{1 - p} (X_{ij} \cdot M_{ij})
\end{align}
During inference ($\text{training} = \text{false}$), the operation is the deterministic identity mapping:
\begin{equation}
Y_{ij} = X_{ij}
\end{equation}

\section{Algorithm Specification}

\begin{algorithm}[H]
\caption{Inverted Dropout Forward Execution}
\label{alg:dropout_inverted}
\begin{algorithmic}[1]
\Require Activation matrix $X \in \mathbb{R}^{B \times D}$, dropout probability $p \in [0, 1)$, mode $\text{training} \in \{\text{true}, \text{false}\}$, optional mask $M_{\text{in}}$.
\Ensure Output matrix $Y \in \mathbb{R}^{B \times D}$, applied mask $M \in \{0, 1\}^{B \times D}$, scale factor $s \in \mathbb{R}_{> 0}$.
\State $B \gets \text{rows}(X)$, $D \gets \text{cols}(X)$
\State Allocate $Y \in \mathbb{R}^{B \times D}$, $M \in \mathbb{R}^{B \times D}$
\If{$\text{training} = \text{false} \lor p = 0$}
    \State $Y \gets X$, $M \gets \mathbf{1}^{B \times D}$, $s \gets 1.0$
    \State \Return $\{Y, M, s\}$
\EndIf
\State $q \gets 1.0 - p$, $s \gets 1.0 / q$
\For{$i \gets 1 \textbf{ to } B$}
    \For{$j \gets 1 \textbf{ to } D$}
        \If{$M_{\text{in}}$ is provided}
            \State $M_{ij} \gets M_{\text{in}}[i][j]$
        \Else
            \State $u \sim \mathcal{U}(0, 1)$
            \State $M_{ij} \gets \mathbf{if } u < q \mathbf{ then } 1.0 \mathbf{ else } 0.0$
        \EndIf
        \State $Y_{ij} \gets X_{ij} \cdot M_{ij} \cdot s$
    \EndFor
\EndFor
\State \Return $\{Y, M, s\}$
\end{algorithmic}
\end{algorithm}

\section{Theoretical Guarantees and Expectation Preservation}

\begin{theorem}[Unbiased Expectation of Inverted Dropout]
For any deterministic input $X \in \mathbb{R}^{B \times D}$ and retention probability $q = 1 - p > 0$:
\begin{equation}
\mathbb{E}_M [Y_{ij}] = X_{ij}
\end{equation}
and the variance induced by the stochastic mask is:
\begin{equation}
\text{Var}_M(Y_{ij}) = \frac{p}{1 - p} X_{ij}^2
\end{equation}
\end{theorem}

\begin{proof}
Because $M_{ij} \sim \text{Bernoulli}(q)$, $\mathbb{E}[M_{ij}] = q = 1 - p$.
Taking the expectation of the inverted dropout activation:
\begin{equation}
\mathbb{E}_M[Y_{ij}] = \mathbb{E}_M \left[ \frac{1}{1-p} X_{ij} M_{ij} \right] = \frac{X_{ij}}{1-p} \mathbb{E}_M[M_{ij}] = \frac{X_{ij}}{1-p} (1-p) = X_{ij}
\end{equation}
For the variance:
\begin{align}
\text{Var}(Y_{ij}) &= \left( \frac{X_{ij}}{1-p} \right)^2 \text{Var}(M_{ij}) = \frac{X_{ij}^2}{(1-p)^2} q(1-q) = \frac{X_{ij}^2}{(1-p)^2} (1-p)p = \frac{p}{1-p} X_{ij}^2
\end{align}
proving exact unbiased expectation preservation.
\end{proof}

\subsection{Limitations}
\begin{itemize}
    \item \textbf{Inference Evaluation Mode Necessity}: Leaving dropout active during production inference introduces unwanted stochastic noise and variance $\mathcal{O}(p/(1-p))$.
\end{itemize}

\section{Complexity Analysis}

\subsection{Time Complexity}
\begin{itemize}
    \item \textbf{Mask Sampling and Scaling}: Requires generating $B \times D$ Bernoulli random bits and performing $B \times D$ element-wise multiplications.
    \item \textbf{Total Time Complexity}: $\Theta(B \cdot D)$ FLOPs.
\end{itemize}

\subsection{Space Complexity}
\begin{itemize}
    \item \textbf{Bitmask Storage}: Storing binary mask for the backward pass requires $\mathcal{O}(B \cdot D)$ bits.
\end{itemize}

\section{Exactness and Error Analysis}

Inverted dropout is algebraically exact. Storing binary masks as packed bitvectors (1 bit per float) in low-level frameworks optimizes memory bandwidth without numerical error.

\section{Worked Numerical Example}

Let $B = 1, D = 4, p = 0.5 \implies q = 0.5, s = 2.0$.
\begin{equation}
X = [[2.0, 4.0, 6.0, 8.0]], \quad M = [[1.0, 0.0, 1.0, 0.0]]
\end{equation}

\begin{enumerate}
    \item \textbf{Scale Calculation}:
    \begin{equation}
    s = \frac{1}{1.0 - 0.5} = \frac{1}{0.5} = 2.0
    \end{equation}
    \item \textbf{Element-wise Application}:
    \begin{align}
    Y_{11} &= 2.0 \times 1.0 \times 2.0 = 4.0 \\
    Y_{12} &= 4.0 \times 0.0 \times 2.0 = 0.0 \\
    Y_{13} &= 6.0 \times 1.0 \times 2.0 = 12.0 \\
    Y_{14} &= 8.0 \times 0.0 \times 2.0 = 0.0
    \end{align}
    \item \textbf{Expectation Verification}:
    \begin{equation}
    \mathbb{E}[Y] = [0.5(4.0) + 0.5(0.0), \dots] = [2.0, 4.0, 6.0, 8.0] = X
    \end{equation}
\end{enumerate}

\section{Numerical Considerations and Edge Cases}

\begin{itemize}
    \item \textbf{Rate $p = 0$}: Identity passthrough ($s = 1.0$).
    \item \textbf{Monte Carlo Dropout}: Running multiple stochastic passes ($T = 50$) at test time enables Bayesian uncertainty estimation.
\end{itemize}

\begin{thebibliography}{9}
\bibitem{srivastava2014dropout} N.~Srivastava, G.~Hinton, A.~Krizhevsky, I.~Sutskever, and R.~Salakhutdinov, ``Dropout: A Simple Way to Prevent Neural Networks from Overfitting,'' \emph{Journal of Machine Learning Research (JMLR)}, vol.~15, no.~56, pp.~1929--1958, 2014.
\bibitem{gal2016dropout} Y.~Gal and Z.~Ghahramani, ``Dropout as a Bayesian Approximation: Representing Model Uncertainty in Deep Learning,'' in \emph{International Conference on Machine Learning (ICML)}, pp.~1050--1059, 2016.
\end{thebibliography}

\end{document}
